\documentclass[sigconf,nonacm]{acmart}

\settopmatter{printacmref=false,printfolios=true}
\renewcommand\footnotetextcopyrightpermission[1]{}
\pagestyle{plain}

\title{Ghostext: an almost perfect stenography via large language model and arithmetic coding}

\author{Anonymous Author(s)}
\affiliation{
  \institution{Anonymous Institution(s)}
  \city{Anonymous City}
  \country{Anonymous Country}
}
\email{anonymous@example.com}

\begin{document}

\begin{abstract}
Ghostext is a linguistic steganography system that combines authenticated encryption with arithmetic-coding-style interval narrowing over large language model next-token distributions. The design goal in this revision is operational reliability under a passive-censor setting: deterministic sender-receiver synchronization, explicit failure semantics, and reproducible configuration binding. We provide an implementation-backed protocol description, configuration-consistent approximation bounds, and initial real-backend baseline results with released artifacts. We explicitly separate measured evidence from open items: approximation diagnostics such as $\widehat{B}_{\mathrm{TV}}$ are defined but not yet fully measured at scale, detector-facing claims remain pending, and active text modification in transit is out of scope in this revision.
\end{abstract}

\maketitle

\section{Introduction}

Steganography and cryptography solve different problems. Cryptography protects message content, while steganography aims to hide that communication is taking place at all. In text channels, this objective is difficult because both lexical quality and token-level statistics can reveal generation artifacts.

Neural Linguistic Steganography (NLS)~\cite{ziegler2019neural} established a practical baseline for arithmetic-coding-based LLM steganography, and OD-Stega~\cite{huang2026odstega} extended this line with constrained-distribution optimization and explicit discussion of tokenization mismatch. Ghostext follows the same generation-based family but focuses on protocol engineering questions that strongly affect deployment: deterministic synchronization, authenticated packet framing, explicit mismatch detection, and fail-closed decode behavior. The long-term aspiration is distribution-preserving steganography, but this draft is intentionally written as a protocol-engineering paper rather than an almost-perfect-security claim.

Ghostext uses the standard three-role setting. Alice takes a prompt, a passphrase, and a secret message. The message is encrypted and then embedded through interval narrowing over model next-token distributions. Bob uses the same protocol-sensitive configuration to decode and decrypt. A passive Censor observes only cover text.

\paragraph{What this paper contributes.}
This paper contributes three concrete elements.

First, it provides an implementation-backed protocol description that maps directly to a companion implementation artifact (anonymized for review), with authenticated packet framing, runtime binding, and fail-closed decode semantics treated as first-class protocol features rather than incidental implementation details.

Second, it introduces a concise formal treatment of approximation sources in the current pipeline: candidate truncation, integer quantization, and retokenization-stability filtering, plus separate discussion of configuration-fingerprint misbinding risk and detector-facing calibration limits for implementation-side diagnostics.

Third, it reports concrete, reproducible end-to-end evidence from the existing implementation: toy-backend bilingual round trips together with a language-stratified real-backend pilot baseline for recoverability, runtime, and efficiency in a template-known regime, while preserving a staged protocol for the remaining approximation-instrumentation, baseline-comparison, prompt-regime, and detector-facing blocks.

\paragraph{Scope boundary.}
This paper reports a pilot operational floor for recoverability/runtime/efficiency, while detector-resistance claims remain intentionally conservative until corresponding measurements are completed. The released 18-trial baseline is therefore evidence of feasibility in a pinned template-known regime, not a settled reliability estimate. Unknown-prompt conditions, larger-sample statistical characterization, aligned NLS/OD-Stega comparison, and active text modification in transit remain outside the measured core of this revision.

\section{Background and Design Motivation}

\subsection{Setting and Notation}

We follow the standard generative steganography setting~\cite{cox2005information,ziegler2019neural}. Let $m$ denote a bitstring to be hidden, and let $y_{1:T}$ denote generated text tokens. At each step $t$, a language model defines a next-token distribution $P_t(\cdot) = P(Y_t \mid y_{<t}, \text{prompt})$ over a vocabulary $\mathcal{V}$. The encoder maps message bits to a token path $y_{1:T}$, and the decoder recovers bits by replaying the same conditional distributions and inverse interval updates.

Ghostext treats $m$ as ciphertext bits, not plaintext bits. This is operationally important because authenticated encryption makes payload bits computationally indistinguishable from random under the shared key assumption. In practical terms, this matches the arithmetic-coding view that expects an almost-uniform source bitstream for efficient mapping.

\subsection{Arithmetic Coding as a Reverse Sampler}

Arithmetic coding is classically introduced as compression~\cite{rissanen1979arithmetic,witten1987arithmetic}. In compression mode, symbol sequences are mapped to compact intervals. In steganography mode, this direction is reversed: a secret bitstream selects a point in an interval domain, and each generation step narrows the interval according to token probabilities until the hidden payload is resolved~\cite{sallee2004model,ziegler2019neural}.

The idealized intuition is straightforward. High-probability tokens occupy larger subintervals and consume fewer bits of uncertainty, while low-probability tokens occupy narrower subintervals and consume more. If encoder and decoder share exactly the same conditional distributions and interval arithmetic, the mapping is invertible.

\subsection{From Ideal Theory to Engineering Reality}

The ideal argument assumes exact distributions, unlimited precision, and perfect tokenization synchronization. Real systems violate all three assumptions.

First, finite precision requires probability quantization and integer interval arithmetic. This introduces approximation error and implementation corner cases that can silently break invertibility if not handled deterministically.

Second, practical candidate policies truncate vocabulary tails for runtime reasons. This alters the effective distribution and can shift the security-capacity tradeoff.

Third, real tokenizer behavior may not guarantee that a detokenized string maps back to one unique token path. OD-Stega~\cite{huang2026odstega} highlights this issue explicitly and treats it as a first-class engineering concern.

These realities motivate Ghostext's design philosophy: prioritize deterministic synchronization and explicit failure over optimistic decoding.

\subsection{Why ``Almost Perfect'' Is a Useful but Bounded Phrase}

NLS~\cite{ziegler2019neural} and OD-Stega~\cite{huang2026odstega} both frame the core tension between payload efficiency and statistical detectability. Ghostext adopts the same framing but with conservative wording.

In this draft, ``almost perfect'' means the following engineering objective under the passive-censor model: preserve model-consistent token behavior as closely as finite quantization and candidate filtering allow, while embedding close to available entropy and preserving deterministic recoverability.

It does \emph{not} mean perfect indistinguishability against all detectors, nor robustness under active editing, nor a proof that finite implementation details cannot leak structure. Those stronger claims require dedicated experiments and, in some cases, separate formal analysis.

\subsection{Positioning Against OD-Style Optimization}

OD-Stega shows that under explicit divergence constraints, temperature-like distribution adjustment can increase utilization in controlled ways~\cite{huang2026odstega}. Ghostext currently does not implement this optimization layer. Instead, it focuses on protocol correctness and reproducibility foundations that are prerequisites for meaningful optimization studies.

This separation is intentional: we first stabilize the encode/decode substrate and its claim-evidence traceability, then add divergence-aware optimization as an incremental extension.

\section{Threat Model}

\subsection{Roles and Shared State}

Ghostext has three roles: Alice (sender), Bob (receiver), and Censor (observer). Alice and Bob share passphrase and protocol-sensitive configuration. Shared configuration includes model/tokenizer identity, prompt, seed, candidate policy, quantization settings, and codec/crypto parameters.

\subsection{Primary Adversary}

The primary adversary is a passive Censor that observes text but does not modify transport content. The Censor does not know the passphrase. Prompt knowledge is treated as an adversarial capability axis: we report our baseline analysis under unknown-prompt assumption and explicitly discuss risks when prompt templates are predictable.

\subsection{Additional Realistic Attack Surfaces}

Beyond pure text inspection, we consider three practical surfaces that require explicit measurement beyond the current baseline.

\textbf{Multi-message accumulation}: detector sees many samples from one sender.

\textbf{Traffic side channels}: detector observes length/timing distributions.

\textbf{Model-family fingerprinting}: detector compares outputs against likely model/prompt families.

Ghostext currently provides defensive engineering hooks for these surfaces (failure taxonomy, retry controls, and configuration binding), but does not yet claim empirical robustness against them.

\subsection{Capability Matrix}

For clarity, Table~\ref{tab:threat-capabilities} summarizes which adversary capabilities are in scope for implemented defenses and which are either covered in baseline measurements or kept as pending measured dimensions.

\begin{table}[t]
\caption{Threat-capability matrix for current Ghostext scope.}
\label{tab:threat-capabilities}
\begin{tabular}{p{0.40\linewidth}p{0.21\linewidth}p{0.31\linewidth}}
\toprule
Adversary capability & Scope status & Current handling \\
\midrule
Passive text observation & in scope & deterministic protocol + fail-closed checks \\
Prompt-template knowledge & measured axis & separate unknown/template-known regimes in evaluation plan \\
Multi-message accumulation & measured axis & explicit detector-aggregation block in staged protocol \\
Traffic timing/length observation & measured axis & retry and latency/length logging with baseline summaries and pending condition expansion \\
Active text modification in transit & out of scope (current version) & no robustness claim; reserved for future extension \\
Runtime compromise of sender/receiver endpoint & out of scope (current version) & no claim; orthogonal endpoint-security problem \\
\bottomrule
\end{tabular}
\end{table}

\subsection{Assumption Boundaries}

Two assumptions deserve explicit caveats. First, sender and receiver need synchronized runtime state, including tokenizer-compatible behavior. Configuration fingerprinting detects drift but does not remove the operational need for configuration hygiene. Second, passphrase secrecy is required for packet confidentiality and integrity; if endpoint secrets are compromised, steganographic concealment alone is insufficient.

These boundaries are not unique to Ghostext, but making them explicit is important for conservative claim wording.

\subsection{Security Objective in This Paper}

The objective of this paper is modest and explicit: provide a deterministic, fail-closed, implementation-auditable steganographic protocol with bounded approximation analysis and a reproducible evaluation framework. Strong detector-evasion claims remain pending until dedicated detector and traffic-side stages are completed.

\section{System Overview}

\begin{figure}[t]
\centering
\fbox{\parbox{0.96\linewidth}{
\textbf{Encode path:}
plaintext $\rightarrow$ UTF-8 bytes $\rightarrow$ authenticated packet $\rightarrow$ split (bootstrap/body) $\rightarrow$ candidate policy + quantization + interval narrowing $\rightarrow$ cover text prefix (+ optional natural tail).

\textbf{Decode path:}
cover text $\rightarrow$ tokenize $\rightarrow$ interval inversion (bootstrap) $\rightarrow$ header authentication + configuration check $\rightarrow$ interval inversion (body) $\rightarrow$ packet authentication $\rightarrow$ plaintext.
}}
\caption{Ghostext end-to-end pipeline. The packet-carrying prefix is deterministic under matched configuration; optional tail does not carry payload.}
\Description{A two-path flow chart showing encode from plaintext to packet to token embedding, and decode from cover text to token inversion to plaintext, with configuration checks and fail-closed gates.}
\label{fig:pipeline}
\end{figure}

Ghostext combines a cryptographic packet layer and an LLM token-selection layer (Figure~\ref{fig:pipeline}). The first protects message confidentiality and integrity. The second provides reversible embedding into generated cover text.

\subsection{Deterministic Synchronization}

Synchronization is a protocol invariant. Ghostext commits runtime settings, backend metadata, and prompt hash into a configuration fingerprint embedded in an authenticated internal header. Decoding proceeds only when this fingerprint matches local runtime state.

\subsection{Two-Stage Embedding Lifecycle}

Embedding is split into bootstrap and body stages. The bootstrap stage carries enough metadata for Bob to authenticate header information and determine expected body length before body decoding. This supports early rejection on mismatch and cleaner failure diagnosis.

\subsection{Natural Tail Handling}

After packet resolution, the encoder may append a short natural tail. Tail tokens do not carry payload and are ignored by the decoder after packet recovery. This isolates surface fluency adjustments from correctness-critical embedding.

\subsection{Fail-Closed Principle}

Mismatch, integrity failure, interval collapse, unstable tokenization paths, and token-budget exhaustion all trigger explicit failure. The decoder never emits partial plaintext after authentication failure.

\section{Protocol and Approximation Model}

\subsection{Authenticated Packet Layer}

Ghostext packet format is
\begin{equation}
\texttt{salt} \parallel \texttt{header\_nonce} \parallel \texttt{sealed\_header} \parallel \texttt{body\_nonce} \parallel \texttt{body\_ciphertext}.
\end{equation}
The internal header stores protocol version, crypto IDs, body length, and configuration fingerprint. Current implementation uses scrypt for key derivation, HKDF-SHA256 for key separation, and ChaCha20-Poly1305 for authenticated encryption.

\subsection{Candidate Policy}

At step $t$, raw logits define base distribution $P_t$ over vocabulary $\mathcal{V}$. Ghostext constructs candidate set $C_t$ via deterministic top-$p$ and max-$k$ truncation, then renormalizes to $\widetilde{P}_t$. If $|C_t|<2$ or step entropy is below threshold, no payload bit is consumed and the top token is emitted.

Retokenization-stability filtering then removes candidates that fail render/re-tokenize path consistency.

\subsection{Integer Quantization}

Let $N_t=|C_t|$ and total frequency budget be $F_{\mathrm{tot}}$ (default $65536$ in current implementation). For each candidate $v\in C_t$:
\begin{equation}
 f_t(v)=1+\left\lfloor \widetilde{P}_t(v)(F_{\mathrm{tot}}-N_t)\right\rfloor,
\end{equation}
with leftover units distributed by fractional remainder and token-id tie-break.

The induced quantized distribution is
\begin{equation}
Q_t(v)=\frac{f_t(v)}{F_{\mathrm{tot}}}.
\end{equation}

\subsection{Interval Update}

Given interval $[L_t,H_t)$ and cumulative boundaries from $Q_t$, selecting token $v$ updates
\begin{align}
L_{t+1}&=L_t+\left\lfloor(H_t-L_t)\,\mathrm{cdf}^-_t(v)\right\rfloor,\\
H_{t+1}&=L_t+\left\lfloor(H_t-L_t)\,\mathrm{cdf}^+_t(v)\right\rfloor.
\end{align}
A segment resolves when $H_t-L_t=1$.

\subsection{Bound for Quantization Error}

For each candidate, deterministic rounding implies
\begin{equation}
\left|Q_t(v)-\widetilde{P}_t(v)\right|\le\frac{1}{F_{\mathrm{tot}}}.
\end{equation}
Hence
\begin{equation}
\|Q_t-\widetilde{P}_t\|_1\le\frac{N_t}{F_{\mathrm{tot}}}, \quad
\mathrm{TV}(Q_t,\widetilde{P}_t)\le\frac{N_t}{2F_{\mathrm{tot}}}.
\end{equation}

\subsection{Bound for Truncation + Quantization + Stability Filtering}

Define truncated tail mass
\begin{equation}
\alpha_t = 1-\sum_{v\in C_t}P_t(v),
\end{equation}
and filtered-out mass under $Q_t$
\begin{equation}
\beta_t = \sum_{v\in C_t\setminus S_t}Q_t(v),
\end{equation}
where $S_t$ is the stability-surviving subset.

Let $R_t$ be $Q_t$ restricted to $S_t$ before renormalization. Then
$\|R_t-Q_t\|_1=\beta_t$ and, because $Q'_t=R_t/(1-\beta_t)$,
$\|Q'_t-R_t\|_1=\beta_t$. Therefore
\begin{equation}
\mathrm{TV}(Q'_t,Q_t)\le\beta_t,
\end{equation}
so the normalization effect after filtering is already covered by the same $\beta_t$ term.

Combining truncation, quantization, and filtering gives
\begin{equation}
\mathrm{TV}(Q'_t,P_t)\le\alpha_t + \frac{N_t}{2F_{\mathrm{tot}}} + \beta_t.
\label{eq:tv_total_bound}
\end{equation}
This bound is implementation-aligned and directly measurable from runtime statistics.

\subsection{Fingerprint Collision Discussion}

Configuration fingerprint uses 64-bit truncation for synchronization mismatch detection. For $n$ independent sessions,
\begin{equation}
\Pr[\mathrm{collision}]\approx\frac{n(n-1)}{2^{65}}.
\end{equation}
For $n=10^6$, this is approximately $2.7\times10^{-8}$. Packet authenticity still depends on AEAD tags; fingerprint collisions do not replace cryptographic integrity checks.

In adversarial settings, a party might search configuration variants that collide in 64 bits. This does not produce valid plaintext release without passphrase-bound AEAD success, but it can degrade diagnostic clarity by making a mismatch appear superficially aligned at the fingerprint layer. For this reason, fingerprint match should be interpreted as a fast guard, while authenticity remains cryptographic.

\subsection{Concrete Encode/Decode Examples}

\begin{table*}[t]
\caption{Two concrete end-to-end examples from current toy-backend implementation.}
\label{tab:examples}
\begin{tabular}{p{0.17\linewidth}p{0.20\linewidth}p{0.40\linewidth}p{0.18\linewidth}}
\toprule
Prompt type & Secret message & Cover text (prefix) & Outcome / token stats \\
\midrule
English walk prompt & \texttt{Meet at 7 PM by the old bridge.} & \texttt{The afternoon feels calm ... people talk softly along the street ...} & decode success; total 455; packet 428; tail 27; bits/token 2.092 \\
Chinese walk prompt (romanized message) & \texttt{Jin wan qi dian zai lao di fang jian.} & \texttt{Jin tian de feng hen qing, jie jiao de kafei dian hai liang zhe deng ...} & decode success; total 424; packet 415; tail 9; bits/token 2.226 \\
\bottomrule
\end{tabular}
\end{table*}

\section{Implementation}

Ghostext is implemented in a companion code artifact (anonymized for review) whose modules map directly to protocol stages: packet crypto and framing, candidate construction with stability checks, integer quantization, interval codec, and end-to-end encode/decode. We provide a deterministic toy backend for protocol-level regression testing and an LLM backend on \texttt{Qwen/Qwen3.5-2B} (third-party GGUF at \texttt{bartowski/Qwen\_Qwen3.5-2B-GGUF}, file \texttt{Qwen\_Qwen3.5-2B-Q4\_K\_S.gguf}; upstream base model \texttt{Qwen/Qwen3.5-2B-Base}) via \texttt{llama-cpp-python} with packaged llama.cpp build \texttt{f49e917}. Released per-run summaries pin Ghostext and paper git heads, hardware/runtime context, tokenizer hash, backend identifier, \texttt{llama-cpp-python} version, build commit, and GGUF metadata. Existing tests cover packet crypto round-trip and mismatch behavior, deterministic quantization, interval invertibility in toy settings, wrong-prompt/seed/text-mutation abort behavior, retokenization stability, bilingual toy round trips, and optional LLM-backend smoke tests. These support correctness and do not imply cross-hardware determinism beyond the pinned local environment.

\section{Security Analysis and Discussion}
\label{sec:analysis}

\subsection{Scope and Goal}

This section analyzes the parts of Ghostext that are already implemented and directly measurable. It is not a complete indistinguishability proof against adaptive steganalysis. The goal is narrower: isolate each approximation source in the current pipeline, express its effect with explicit quantities, and define what remains to be measured beyond the released pilot baseline.

\subsection{Distribution Pipeline Decomposition}

At token step $t$, the implementation transforms the base next-token distribution $P_t$ through three deterministic operators:
\begin{equation}
P_t \xrightarrow{\text{truncate}} \widetilde{P}_t \xrightarrow{\text{quantize}} Q_t \xrightarrow{\text{stability filter}} Q'_t.
\end{equation}

The truncation operator keeps a top-$p$/max-$k$ subset and removes tail mass. Quantization maps real probabilities to integer frequencies with total budget $F_{\mathrm{tot}}$. Stability filtering removes candidates that fail render/re-tokenize consistency and then renormalizes.

This decomposition is useful because each term is observable in logs and can be ablated independently.

Configuration note: the released real-backend baseline uses $F_{\mathrm{tot}}=4096$ and $h_{\min}=0.0$, while implementation defaults are $F_{\mathrm{tot}}=65536$ and $h_{\min}=1.0$. All numeric examples below are labeled explicitly, and the released efficiency numbers should be interpreted as an ungated initial operating point rather than a recommended deployment setting.

\subsection{Per-Step Total-Variation Bound}

Section~\ref{sec:protocol-details} defines
\begin{equation}
\alpha_t = 1-\sum_{v\in C_t} P_t(v),
\end{equation}
\begin{equation}
\beta_t = \sum_{v\in C_t\setminus S_t} Q_t(v),
\end{equation}
and provides
\begin{equation}
\mathrm{TV}(Q'_t,P_t)\le\alpha_t + \frac{N_t}{2F_{\mathrm{tot}}} + \beta_t,
\end{equation}
where $N_t=|C_t|$.

The three addends have clear meanings.
$\alpha_t$ is policy truncation mass,
$N_t/(2F_{\mathrm{tot}})$ is finite-frequency quantization distortion,
and $\beta_t$ is additional mass removed by tokenization-stability filtering, including the normalization effect after filtering.

\subsection{Quantization Bias and a Conservative KL Upper Bound}

To answer the reviewer concern on quantization bias, define
$\delta_t(v)=Q_t(v)-\widetilde{P}_t(v)$.
By construction, $\sum_v\delta_t(v)=0$ and $|\delta_t(v)|\le 1/F_{\mathrm{tot}}$.
Because every quantized candidate gets at least frequency one, $Q_t(v)\ge 1/F_{\mathrm{tot}}$.

Using $\widetilde{P}_t(v)=Q_t(v)-\delta_t(v)$:
\begin{align}
D_{\mathrm{KL}}(\widetilde{P}_t\|Q_t)
&=\sum_{v:\widetilde{P}_t(v)>0} \widetilde{P}_t(v)\log\!\left(\frac{\widetilde{P}_t(v)}{Q_t(v)}\right).
\end{align}
Applying the standard $D_{\mathrm{KL}}(P\|Q)\le \chi^2(P,Q)$ upper bound~\cite{tsybakov2009introduction} gives
\begin{align}
D_{\mathrm{KL}}(\widetilde{P}_t\|Q_t)
&\le \chi^2(\widetilde{P}_t,Q_t)
= \sum_{v:\widetilde{P}_t(v)>0} \frac{(\widetilde{P}_t(v)-Q_t(v))^2}{Q_t(v)}
= \sum_{v:\widetilde{P}_t(v)>0} \frac{\delta_t(v)^2}{Q_t(v)}.
\end{align}
Then, using $Q_t(v)\ge 1/F_{\mathrm{tot}}$ and $|\delta_t(v)|\le 1/F_{\mathrm{tot}}$,
\begin{align}
D_{\mathrm{KL}}(\widetilde{P}_t\|Q_t)
&\le \sum_v \frac{(1/F_{\mathrm{tot}})^2}{1/F_{\mathrm{tot}}}
\le \frac{N_t}{F_{\mathrm{tot}}}\quad\text{(nats)}.
\end{align}

The resulting bound is loose but explicit. Under implementation defaults ($N_t\le16$, $F_{\mathrm{tot}}=65536$), it is at most $16/65536 \approx 2.44\times10^{-4}$ nats for quantization alone. Under the released baseline setting ($N_t\le64$, $F_{\mathrm{tot}}=4096$), the corresponding worst-case ceiling is $64/4096 = 1.56\times10^{-2}$ nats. Truncation and stability filtering can dominate this term in practice, so all terms should be logged together.

\subsection{Entropy-Threshold Gating and Capacity}

Embedding is allowed only when both conditions hold: $N_t\ge2$ and per-step entropy is at least $h_{\min}$.
The implementation default is $h_{\min}=1.0$ bit, while the released real-backend baseline sets $h_{\min}=0.0$ to avoid entropy-threshold gating during this initial recoverability/runtime run.

Let $I_t\in\{0,1\}$ indicate whether step $t$ allows embedding, and let $\Delta b_t$ denote interval-resolved bits at that step.
An implementation-level effective payload rate can be summarized as
\begin{equation}
\eta_{\mathrm{pkt}}=\frac{\sum_t I_t\,\Delta b_t}{\text{packet-carrying tokens}}.
\end{equation}

The threshold therefore trades capacity for synchronization safety.
Lower $h_{\min}$ increases opportunities to embed, but also increases exposure to low-entropy stalls and retry events. Accordingly, the released $2.478$ bits/token figure should be read as an ungated initial baseline under $h_{\min}=0.0$, not as a safe-default deployment recommendation.
The practical next step is therefore not to headline this single number, but to run the same condition grid at least over $h_{\min}\in\{0.0,0.5,1.0\}$ before treating any efficiency figure as representative of a conservative deployment envelope.

\subsection{Sequence-Level Accumulation and Diagnostics}

For full-sequence distributions, chain decomposition gives
\begin{equation}
D_{\mathrm{KL}}(Q_{1:T}\|P_{1:T})
=\sum_{t=1}^{T}\mathbb{E}_{Q_{<t}}\left[D_{\mathrm{KL}}(Q_t(\cdot\mid Y_{<t})\|P_t(\cdot\mid Y_{<t}))\right].
\end{equation}

In Ghostext, contexts depend on prior stego choices, so a tight analytic bound is difficult without stronger assumptions. The practical path is to log per-step terms and compute empirical diagnostics per message:
\begin{equation}
\widehat{B}_{\mathrm{TV}}=\sum_{t=1}^{T}\left(\alpha_t + \frac{N_t}{2F_{\mathrm{tot}}} + \beta_t\right).
\label{eq:bhat-tv}
\end{equation}

If $P_{1:T}$ denotes the matched-distribution natural-cover process, then sequence-level distinguishability is governed by $\mathrm{TV}(Q_{1:T},P_{1:T})$. With $\mathrm{Adv}_{\mathrm{censor}} := |2\Pr[\mathrm{correct}] - 1|$ for a binary hypothesis test between one sample from $Q_{1:T}$ and one sample from $P_{1:T}$, an idealized censor satisfies
\begin{equation}
\mathrm{Adv}_{\mathrm{censor}}\le \mathrm{TV}(Q_{1:T},P_{1:T}) \le \sqrt{\tfrac{1}{2}D_{\mathrm{KL}}(Q_{1:T}\|P_{1:T})},
\end{equation}
where the first inequality is the standard testing bound and the second is Pinsker's inequality~\cite{tsybakov2009introduction}. The steganographic meaning of that bridge is then interpreted through frameworks such as Cachin~\cite{cachin1998information} and Hopper et al.~\cite{hopper2002provably}. That sequence-level bridge is the quantity a censor-facing claim would ultimately need.

$\widehat{B}_{\mathrm{TV}}$ is therefore an audit signal rather than a direct detector-risk guarantee. In this revision we use it conservatively in two ways: a large value is direct evidence that approximation terms are not negligible, while a small value is only suggestive and still requires detector-side validation. Its practical usefulness is therefore calibration-dependent: if a later regime shows small $\widehat{B}_{\mathrm{TV}}$ but high detector AUC, the metric should be downgraded to a local implementation-health diagnostic rather than treated as a censor-risk proxy. For that reason, the staged detector block treats correlation analysis between $\widehat{B}_{\mathrm{TV}}$ and detector scores as a required output, not an optional add-on. The full real-backend measurement block remains pending (Stage 3).

\subsection{Relation to Provable-Security Limits}

A useful limit case is $\alpha_t\to0$ (no truncation), $\beta_t\to0$ (no stability filtering loss), and $F_{\mathrm{tot}}\to\infty$ (no quantization loss), under exact sender-receiver replay of the same next-token distribution. In this limit, the implementation approaches ideal arithmetic-coding steganography over the model distribution at each step.

This limit clarifies where Ghostext sits relative to Cachin-style distribution matching and Hopper-style stego security formulations: our current system is an approximation to that ideal regime, with approximation terms made explicit by $\alpha_t$, $N_t/(2F_{\mathrm{tot}})$, and $\beta_t$. The present paper does not claim that current finite-precision, truncated, stability-filtered operation already satisfies a full provable-security theorem.

\subsection{Fingerprint Collision and Misbinding}

The synchronization fingerprint is a truncated 64-bit SHA-256 digest over runtime configuration, backend metadata, and prompt hash. For $n$ honest sessions,
\begin{equation}
\Pr[\text{collision}]\approx\frac{n(n-1)}{2^{65}}.
\end{equation}
For $n=10^6$, this is roughly $2.7\times10^{-8}$.

A collision is not sufficient for successful forgery. Decoding still requires authenticated packet validity under the shared passphrase and protocol checks. The fingerprint should be interpreted as a fast mismatch detector, not as a stand-alone authenticity primitive.
Under an idealized uniform-target approximation, $2^{40}$ adaptive collision attempts against a fixed 64-bit target still yield only about $2^{-24}$ success probability, which is small but not a margin worth presenting as future-proof when the byte overhead of widening is negligible. The fingerprint is also a Bob-local synchronization check rather than a public plaintext header field in the passive model of this paper. The 64-bit width is therefore a deployment tradeoff in the current artifact, not a claim that 64 bits is universally optimal. Given the negligible byte overhead of widening this field, 128-bit fingerprinting is the conservative future direction unless compatibility with already released artifacts must be preserved.

\subsection{Retry-Induced Traffic Leakage Model}

Retries improve robustness in difficult regions but can change externally observable traffic.
Let $R$ be attempt count, $L$ output length, and $T_{\mathrm{lat}}$ latency.
Observed marginals become mixtures:
\begin{equation}
\Pr[L=\ell]=\sum_r \Pr[R=r]\Pr[L=\ell\mid R=r],
\end{equation}
\begin{equation}
\Pr[T_{\mathrm{lat}}\le t]=\sum_r \Pr[R=r]\Pr[T_{\mathrm{lat}}\le t\mid R=r].
\end{equation}

Importantly, $R$ is sender-local: failed attempts are discarded and Bob only decodes the final transmitted cover. So $R$ does not alter decode replay correctness, but it can affect sender-side timing/energy and external traffic observables.

Even when text-only quality is strong, these mixture shifts may be detectable by a traffic observer. In the released 18-trial pilot, \texttt{attempts\_used} has mean $1.778$ and maximum $4$: 11 runs finished in one attempt, 2 in two attempts, 3 in three attempts, and 2 in four attempts. The English subset is mostly single-attempt (7 of 9), while the Chinese subset spreads across one, two, three, and four attempts with counts 4, 1, 2, and 2. By contrast, a natural-cover generator without a retry loop has degenerate $R=1$ unless dummy retries are added deliberately. Under the current policy, retry count alone is therefore a trivially separating feature against any zero-retry natural-cover baseline; no learned detector is required to exploit that difference. This already makes the side-channel concern concrete in the released pilot, and it also explains why the released p99 encode latency is driven by a few high-retry cases rather than a stabilized tail law. Evaluation should therefore report retry-conditioned and unconditional length/latency distributions, and we keep this as a pending detector/traffic block beyond the current baseline.

\subsection{Prompt-Knowledge and Multi-Message Regimes}

The released real-backend evidence is a template-known operational floor: prompts are fixed public templates repeated across cases. For broader deployment claims, two stronger regimes still need direct measurement:
(1) private-prompt or unknown-prompt conditions, where the censor does not know the sender's exact prompt,
and (2) sample-accumulation, where many messages from one sender are available.

If a detector score is $s(x)$ per message and the censor aggregates
$\bar{s}_m = \frac{1}{m}\sum_{i=1}^{m}s(x_i)$,
small per-message effects can become detectable as $m$ grows.
This is why single-message and multi-message tests should both be first-class reporting targets.

\subsection{Synchronization Assumptions in Deployment}

Current evidence is strongest for local, pinned deployments where both parties control model files, tokenizer behavior, backend build, and runtime policy. The implementation binds backend metadata and tokenizer hash into the configuration fingerprint, which is effective for detecting drift in this deployment mode.

When using third-party hosted APIs, exact backend-version observability may be limited. In that environment, strict deterministic replay can be unavailable by construction, and Ghostext should be treated as out-of-scope unless an API provider exposes stable and verifiable model/tokenizer/version identifiers sufficient for replay.

\subsection{Cross-Hardware and Numeric Determinism Caveat}

The current fingerprint includes backend metadata exposed by the backend implementation, but it does not include a full hardware fingerprint (for example, GPU model and driver stack). In local llama.cpp CPU mode this is less exposed to GPU-kernel nondeterminism than many hosted GPU pipelines, but floating-point and platform differences can still perturb logits near candidate cut boundaries.

The practical implication is conservative: cross-hardware deployment should be validated empirically before operational use, and hardware/runtime context should be logged as a reproducibility field in evaluation artifacts. The current paper therefore claims deterministic replay only within a pinned local environment, not as a universal cross-hardware property.

\subsection{Natural Tail Policy Clarification}

Current natural-tail tokens are sampled from the same constrained candidate mechanism used by packet-carrying tokens, then excluded from decode once packet recovery is complete. This avoids a second sampling pathway but may preserve policy artifacts in tail regions.

The correct interpretation is therefore operational, not security-complete:
in this version, natural tails mainly extend text to a more complete ending, but they are not a distinct detector-evasion mechanism. They should be measured separately from packet-carrying prefixes in both text and traffic analyses.

\subsection{Logging Schema for Measurable Analysis}

\begin{table}[t]
\caption{Per-step and per-message diagnostics and metadata to log for analysis, ablation, and reproducibility.}
\label{tab:analysis-logs}
\begin{tabular}{p{0.25\linewidth}p{0.68\linewidth}}
\toprule
Field & Role in analysis \\
\midrule
$\alpha_t$ & Truncation mass removed from base model distribution at step $t$. \\
$N_t$ & Candidate count after truncation and before quantization. \\
$F_{\mathrm{tot}}$ & Integer frequency budget used by quantizer. \\
$\beta_t$ & Quantized mass removed by retokenization-stability filtering. \\
$H_t$ & Candidate entropy used for min-entropy gating and low-entropy window checks. \\
$I_t$ & Embedding-enabled indicator at step $t$ (1 if payload bits may be consumed). \\
$R$ & Retry count per message; needed for latency/length mixture analysis. \\
Failure class & One of explicit fail-closed classes (integrity, mismatch, stall, budget, etc.). \\
Model metadata & Model family/version and tokenizer hash used in generation. \\
Policy config & top-$p$, max-$k$, min-entropy threshold, total frequency budget, retry limits. \\
Prompt regime & Prompt family label and whether template-known or private-prompt setting. \\
Runtime context & Hardware class, backend build info, and software versions. \\
\bottomrule
\end{tabular}
\end{table}

\subsection{Current Claims Only}

At this stage, supported claims are restricted to deterministic decode under matched runtime state, explicit mismatch detection, fail-closed behavior, and measurable decomposition of approximation sources.

Claims about broad detector resistance and cross-family generalization remain pending dedicated detector and cross-family experiments. Active text modification in transit remains out of scope in this revision, and deterministic claims should be read as scoped to the pinned local environment used for the released artifact.

\section{Evaluation}

We keep the evaluation narrow and tied to the released implementation. Instead of relying on a single smoke-style benchmark, we report a fixed-grid experiment driven by the released Ghostext evaluation script \nolinkurl{run_prompt_message_grid.py}. The grid crosses ten public prompts with ten public messages for 100 end-to-end encode/decode trials under the same pinned local backend from \S\ref{sec:open-science}. The prompt set contains five English and five Chinese templates. The message set contains four short, four medium, and two long plaintexts spanning both languages. All trials use the released defaults. We keep seed $7$ throughout. The decoding settings remain \texttt{top\_p}=0.995, \texttt{max\_candidates}=64, \texttt{min\_entropy\_bits}=0.0, and \texttt{total\_frequency}=4096. The built-in low-entropy retry logic stays enabled throughout. A fixed local passphrase is reused across the grid for reproducibility, but only the policy and not the secret itself is released.

\subsection{Fixed 10x10 Prompt-Message Grid}

\begin{table}[t]
\caption{Fixed 10x10 prompt-message grid on the pinned local \texttt{llama-cpp} backend.}
\label{tab:grid-benchmark}
\footnotesize
\centering
\setlength{\tabcolsep}{2pt}
\begin{tabular}{@{}lrrrrrr@{}}
\toprule
Slice & Trials & Success & Enc med (s) & Dec med (s) & Mean bpt & Attempts mean \\
\midrule
Overall & 100 & 100/100 & 37.74 & 34.47 & 2.619 & 1.40 \\
EN prompts & 50 & 50/50 & 35.04 & 34.21 & 2.686 & 1.18 \\
ZH prompts & 50 & 50/50 & 38.90 & 37.09 & 2.552 & 1.62 \\
Short msgs & 40 & 40/40 & 27.90 & 26.45 & 2.645 & 1.08 \\
Medium msgs & 40 & 40/40 & 42.31 & 37.43 & 2.594 & 1.48 \\
Long msgs & 20 & 20/20 & 60.22 & 58.11 & 2.615 & 1.90 \\
\bottomrule
\end{tabular}
\end{table}

Table~\ref{tab:grid-benchmark} shows exact round-trip recovery for all 100 trials, with decode success rate $1.000$ and Wilson 95\% confidence interval $[0.963, 1.000]$. No decode mismatch or fail-open behavior appears in the released run log. Overall median encode latency is $37.74$\,s and median decode latency is $34.47$\,s. Mean encode payload density is $2.619$ bits/token (median $2.677$), while mean packet length is $167.1$ bytes.

The main performance caveat is not correctness but tail latency. Encode latency has mean $50.56$\,s, $p90=74.69$\,s, $p99=196.55$\,s, and maximum $291.69$\,s, whereas decode latency is tighter with mean $38.02$\,s, $p90=58.07$\,s, and $p99=72.22$\,s. Retry behavior explains much of that spread, because some local contexts trigger low-entropy or no-stable-candidate retries. English-prompt cases are modestly faster and denser than Chinese-prompt cases, while message length mainly affects latency rather than mean bits/token.

These results support one narrow claim: within one pinned local environment and one fixed prompt-message grid, Ghostext provides deterministic, fail-closed end-to-end recoverability with transparent overhead measurements. They also give a more realistic view of runtime than the earlier smoke benchmark, because they expose retry-driven tail behavior across multiple prompt and payload regimes. They do not establish detector resistance, direct distribution-distance bounds, cross-hardware determinism, or robustness to active text modification.

\section{Related Work}

NLS~\cite{ziegler2019neural} established arithmetic-coding neural steganography as a practical baseline. Later work explored adaptive and near-imperceptible variants~\cite{shen2020near,dai2019towards} and stronger formal constructions~\cite{zhang2021provably,ding2023discop,dewitt2024perfect}.

OD-Stega~\cite{huang2026odstega} is most directly related to Ghostext in terms of constrained-distribution perspective and tokenization mismatch emphasis.

Foundational information-theoretic and provable-security frameworks~\cite{cachin1998information,hopper2002provably} clarify what a complete steganographic security claim requires.

Ghostext is positioned as a systems-security implementation contribution: deterministic protocol behavior, synchronization binding, explicit failure semantics, and measurable approximation decomposition in an end-to-end codebase.

\subsection{Method Lineage and Security-Claim Styles}

The modern language-steganography line can be read along two axes. One axis is embedding mechanism, where arithmetic-coding-style approaches dominate practical neural systems~\cite{ziegler2019neural,shen2020near}. The second axis is claim style: some works primarily report empirical imperceptibility and utility trends, while others aim for formal indistinguishability guarantees under explicit assumptions~\cite{zhang2021provably,ding2023discop,dewitt2024perfect}.

Ghostext sits closer to the engineering side of this map. It keeps the arithmetic-coding family mechanism, but focuses on deterministic replay, fail-closed runtime semantics, and configuration binding under real implementation constraints. This is a narrower contribution than a new steganographic primitive, and the paper intentionally describes it in that scope.

\subsection{Practical Friction Points in Prior Systems}

Earlier neural approaches generally assume that sender and receiver can replay equivalent model distributions. In deployed systems, this assumption is stressed by tokenizer behavior, backend version drift, and finite-precision codec choices. OD-Stega explicitly discusses tokenizer mismatch risks~\cite{huang2026odstega}, which directly motivates Ghostext's retokenization-stability filtering and explicit mismatch checks.

Another practical gap is failure semantics. Many prior papers report aggregate performance but provide less structured runtime failure taxonomy for deployment troubleshooting. Ghostext emphasizes class-labeled failures because operational debugging and security interpretation often depend on distinguishing integrity failure, configuration mismatch, synchronization drift, and low-entropy stalls.

Kaptchuk et al.~\cite{kaptchuk2021language} also emphasize that language channels can leak secret-dependent structure beyond the intended communication path, reinforcing the need for explicit protocol boundaries and conservative claim wording.

A closely related neighboring line is LLM watermarking, which also manipulates generation distributions under detectability constraints~\cite{kirchenbauer2023watermark}. Although watermarking and steganography have different goals, both depend on distribution-level control and detector assumptions, so cross-citation is useful for threat interpretation.

On the detector side, practical evaluation commonly combines lightweight text classifiers and stronger model-based detectors; examples include fastText-style baselines~\cite{joulin2017fasttext} and curvature-based zero-shot detectors such as DetectGPT~\cite{mitchell2023detectgpt}. These references motivate our detector-block design in the evaluation protocol.

\subsection{Relation to Formal Frameworks}

Cachin's information-theoretic model~\cite{cachin1998information} and Hopper et al.'s provable framework~\cite{hopper2002provably} remain the conceptual baseline for strong security claims. Recent generative-text works with formal guarantees~\cite{zhang2021provably,ding2023discop,dewitt2024perfect} move closer to these standards under specific model-access assumptions.

Ghostext does not claim to replace those frameworks. Instead, it contributes implementation-level decomposition terms and an auditable evidence pipeline with initial baseline results, intended to bridge engineering practice and stronger claims. The design choice is to make approximation sources observable first, then upgrade claims only after corresponding measurements are available.

\subsection{Scope-Constrained Novelty Statement}

Given this landscape, Ghostext claims novelty at the systems-integration level: combining authenticated packet framing, runtime configuration binding, retokenization-aware candidate control, and explicit failure observability in one deterministic pipeline. Whether this combination yields significant detector-facing gains remains an open empirical question for the remaining detector-focused stages.

\section{Limitations and Conclusion}

\paragraph{Limitations.} Ghostext targets passive-text observation and does not provide empirical detector-resistance evidence. LLM-backend evidence remains a pilot 18-trial template-known baseline and a six-case aligned slice; cross-family generalization is not yet established. The $2.448$ bits/token figure is from an ungated baseline and not a deployment recommendation. Deterministic replay is evidenced only within a pinned local environment. Active text modification in transit is out of scope.

\paragraph{Conclusion.} Ghostext contributes a near-distribution-preserving, implementation-auditable LLM steganography protocol with authenticated packetization and synchronization binding. Relative to the same-policy natural-cover baseline, only two implementation-level effects remain: integer quantization and retokenization filtering. The 18-trial LLM-backend baseline is a pilot measurement, and the six-case same-backend slice separates payload-normalization effects from broader detector claims. Completing detector-facing, traffic-side, and cross-family evaluation under the same artifact discipline is the next milestone.


\bibliographystyle{ACM-Reference-Format}
\bibliography{refs}

\appendix
\section{Claim-to-Evidence Mapping (Condensed)}

This appendix is a compact in-paper mirror of the full traceability note in the artifact package. Each claim is scoped to implementation and tests available at draft time. Relative paths refer to the anonymized companion artifact layout.

\begin{table}[h]
\caption{Condensed module-to-evidence traceability map; see Section~\ref{sec:related-positioning} for the matching scope-constrained novelty statement.}
\label{tab:appendix-claim-evidence}
\begin{tabular}{p{0.28\linewidth}p{0.24\linewidth}p{0.39\linewidth}}
\toprule
Claim area & Test / artifact group & Released evidence path or anchor \\
\midrule
Packet crypto and framing (\texttt{crypto.py}, \texttt{packet.py}) & \texttt{test\_crypto.py}, \texttt{test\_packet.py} & companion artifact modules + toy round-trip checks \\
Configuration fingerprint fail-closed behavior (\texttt{decoder.py}) & \texttt{test\_failures.py} & mismatch/fail-closed cases \\
Candidate policy and retokenization stability (\texttt{candidate\_policy.py}) & \texttt{test\_tokenization\_stability.py}, \texttt{test\_failures.py} & candidate/stability cases \\
Deterministic integer quantization (\texttt{quantization.py}) & \texttt{test\_quantization.py} & quantization regression cases \\
Interval codec invertibility (\texttt{codec.py}) & \texttt{test\_codec\_toy.py} & deterministic toy round-trip cases \\
Retry/failure taxonomy in real runs & released baseline logs & \texttt{results/real-backend-baseline-r1/} and \texttt{r2/} \\
Bilingual operational floor and merged metrics & released summaries + smoke coverage & \texttt{results/real-backend-baseline-summary-merged.\{json,md\}} and \texttt{test\_llama\_cpp\_integration.py} \\
\bottomrule
\end{tabular}
\end{table}


\end{document}
